\section{Introduction}
\label{sec:introduction}

\subsection{Two readings, one program}

A smart contract that hosts a sealed-bid auction is one program with
two readings.  Read as an Ethereum contract, it is a state machine
exposing transactions: \TT{commit(\textit{hash})},
\TT{reveal(\textit{bid})}, \TT{settle()}.  Read as a game, it is an
extensive form with players, information sets, and a payoff
function.  These two readings are seldom kept in sync.  The
contract is written, audited, deployed, and run as a state
machine; the game is sketched, sometimes proved correct on paper,
almost never mechanically connected to the bytecode the validators
actually execute.

The gap is where strategic surprises live.  A sealed-bid auction
whose commit phase happens to publish enough information to be
front-run.  An escrow whose timeout branch creates a griefing
incentive that no equilibrium analysis predicted because the
model omitted the timeout.  A voting scheme whose Solidity
implementation reveals the tally one block earlier than the
analysis assumed.

Vegas is a small programming language and compiler aimed at this
specific gap.  A Vegas program is a finite multi-party protocol
written in the natural style of a referee: who joins, who plays
what, what is hidden and revealed when, and how the pool is split
at the end.  Programs are short --- the running examples in this
report are a dozen lines each --- and look more like the textbook
description of a game than like a smart-contract framework.  The
compiler reads the program, infers a dependency relation among
its actions, and emits a family of artifacts that all denote the
same protocol from different viewpoints: a Solidity contract for
deployment, an extensive-form game for equilibrium analysis, an
SMT-LIB encoding for reachability queries, and several others.

The premise is not that strategic correctness is proved
automatically.  The premise is that game-theoretic reasoning
about the program transfers, modulo the stated execution model,
to the deployed contract: if an analyst (or a tool) concludes that
truthful bidding is dominant in the source-language Vickrey
specification, that conclusion carries over to the Solidity
contract running on chain.  EFG generation is one channel for
that reasoning; hand analysis on the source is another;
mechanical proof in a companion Lean library is a third.  Each is
made possible by the same property: the compiler structures the
source so that the program \emph{is} a game with a public history,
information sets given by the visibility discipline, and a
payoff function over the terminal state, and every backend is a
lowering of that game.

This report describes how that compilation is organized, what
the language can and cannot express, how the central
intermediate representation --- the \emph{event graph} ---
captures the strategic structure of a protocol, and how the
per-backend lowerings are constrained so that the artifacts agree.

\subsection{What Vegas is and is not}

Vegas is deliberately narrow.  The aim is to make the round trip from
specification to deployed contract to game-theoretic analysis short
and reliable for a useful class of protocols, not to be a general
purpose smart-contract language.  Concretely, in scope are:

\begin{itemize}\itemsep1pt
\item finite, terminating multi-party protocols with a single
  pool of locked funds;
\item bounded types (booleans, finite ranges, enumerations);
\item deterministic guards (\TT{where} clauses) over prior public
  values and over the value being written;
\item public randomness via an explicit \texttt{random} construct;
\item commit-reveal patterns, with automatic insertion when the
  compiler detects simultaneous public moves that would otherwise be
  front-runnable;
\item explicit operational handling of non-responsive participants
  (timeouts and bail-outs, with handler clauses chosen by the
  programmer);
\item terminal allocation of the pool by a \TT{withdraw} clause.
\end{itemize}

Out of scope, on purpose:

\begin{itemize}\itemsep1pt
\item unbounded data and loops --- the event graph must be finite;
\item ERC-20/ERC-777 token flows and multi-token games (planned, not
  in the current language);
\item interactions with arbitrary other contracts (the protocol is a
  closed system; cross-contract MEV is parameter, not modelled);
\item private randomness from outside oracles;
\item dynamic player counts (the role set is fixed at compile time).
\end{itemize}

These restrictions are real costs.  They are also what make the
multi-backend story feasible: every artifact is computed by
exhaustive enumeration or simple recursion over a finite graph, with
no need to approximate or to choose a sound-but-incomplete encoding.
A Gambit extensive-form game can be built directly from the
EventGraph because the graph is finite; an SMT-LIB encoding of
``can role $R$ reach a configuration where they get $\geq x$ wei?''
is quantifier-free over a finite set of action witnesses; a Lean or
Coq witness-record encoding terminates by construction.  We expect
the language to grow --- variable deposits, tokens, bounded
iteration --- and we have tried to make the IR robust enough to
accommodate those extensions.  The current paper documents the
foundations.

\subsection{A first example}
\label{sec:intro:first-example}

The following Vegas program is the complete specification of a
distribution-semantics Monty Hall game.  Two roles, Host and Guest,
each lock 20 wei of deposit into a 40 wei pool.  The Host hides the
car behind one of three doors; the Guest picks a door; the Host
reveals a goat door; the Guest decides whether to switch; the Host
reveals the car; payouts are computed from the public transcript.
If a role fails to act in time, the \texttt{or split} handler
forfeits their deposit to the surviving party.

\begin{lstlisting}[style=vg]
type door = {0, 1, 2}

game main() {
  join Host()  $ 20;
  join Guest() $ 20;
  commit or split Host(car: door);
  yield  or split Guest(d: door);
  yield  or split Host(goat: door) where Host.goat != Guest.d;
  yield  or split Guest(switch: bool);
  reveal Host(car: door) where Host.goat != Host.car;
  withdraw { Guest -> ((Guest.d != Host.car) <-> Guest.switch) ? 40 : 0;
             Host  -> ((Guest.d != Host.car) <-> Guest.switch) ? 0  : 40 }
}
\end{lstlisting}

From this twelve-line source Vegas produces every artifact listed
in the abstract: a roughly two-hundred-line Solidity contract that
enforces the protocol step by step (with cryptographic commitments
for the Host's hidden door and a fall-back to the \TT{or split}
clauses when a role times out); an extensive-form game in
Gambit's format, containing the move tree and the information
sets induced by the dependency relation; an SMT-LIB
witness-existence encoding; and a Coq/Lean witness-record
encoding (the Diamond DAG of \S\ref{sec:other:diamond}).  All are
produced by the current compiler and exercised by the
golden-master test suite.  The Solidity output is deployable to a
local Anvil node; fed to Gambit, the EFG returns the
``always switch'' equilibrium for Monty Hall.  Nothing in the
twelve-line source mentions Ethereum, cryptographic commitments,
timeouts, gas, or game trees.

The example is small enough to verify by inspection that the two
artifacts agree on what they say is hidden, what is public, when
each player observes what, and which deviations are punished.  At
larger scale the agreement has to be argued, not inspected --- which
is what motivates the companion Lean formalization VegasCore
discussed in \S\ref{sec:vegascore}.

\subsection{Why an event graph}
\label{sec:intro:why-graph}

A natural alternative would be to compile the surface syntax
directly to a state machine for Solidity and directly to a tree
for Gambit, twice, with the two backends owning their own
semantics.  That is also what an experienced human writes by hand
when they want both views of the same protocol.  Vegas does not
do this for three reasons.

First, the textual order of the Vegas source does \emph{not} encode
the strategic structure.  In the Monty Hall example above, the line
\TT{yield Guest(d: door)} appears below the Host's hidden
\TT{commit}.  But the Guest cannot read the Host's hidden value, and
the Host's commit has no data dependency on the Guest; the two
actions are concurrent in the strategic sense, regardless of the
order they happen to be typed in.  Compiling to a tree directly
would force an arbitrary ordering choice, which then has to be
``undone'' with information-set equivalences.  Inferring
concurrency from data dependencies once, in the IR, and letting
every backend read it off the same graph, is simpler.

Second, the operational behaviour of a Solidity contract is more
naturally described by a dependency relation than by a global step
counter.  The Solidity backend emits a \TT{depends} modifier per
action, the modifier gating execution on the completion flags of
the action's predecessors in the event graph; concurrent actions
can arrive in any order.  This matches the on-chain reality (the
transaction ordering is chosen by builders, not by the contract)
without the contract having to encode a totalisation.

Third, the rewriting passes the compiler performs --- in particular,
automatic commit-reveal insertion --- are dependency-graph rewrites.
They are easier to state and easier to test on a graph than on the
syntax tree.

So Vegas uses a single intermediate representation, the
\emph{EventGraph}, between the surface language and every backend.
Each backend reads the graph; no backend reaches behind the graph
to the surface syntax.  Section~\ref{sec:eventgraph} defines the
EventGraph in detail; the rest of the report uses it pervasively.

\subsection{Contributions and outline}
\label{sec:intro:contributions}

The contributions of this report are:

\begin{enumerate}\itemsep2pt
\item \emph{A language design} (\S\ref{sec:language}) that makes it
  feasible to describe a finite on-chain protocol in fewer than
  twenty lines and that compiles, without further specification,
  to both deployable code and a game-theoretic model.

\item \emph{The EventGraph IR} (\S\ref{sec:eventgraph}) and its
  associated invariants (acyclicity, commit-reveal ordering,
  visibility on guard reads), capturing concurrency, observability,
  and the protocol's information-set structure in one finite
  object.

\item \emph{Automatic commit-reveal insertion}
  (\S\ref{sec:commit-reveal}): a graph-rewriting pass that detects
  concurrent public writes and replaces each one with a
  commit-then-reveal pair, with predecessors set up so that no
  player can observe another player's reveal before committing
  themselves.  This addresses one specific subclass of MEV ---
  observation-before-action on simultaneous moves within the
  protocol --- by making the mitigation a structural property
  of the IR rather than an implementation discipline.

\item \emph{A family of backends} (\S\S\ref{sec:solidity}--\ref{sec:other-backends})
  that each consume the same EventGraph: Solidity and Vyper for
  EVM deployment with \TT{depends}-modifier scheduling and a
  timeout-driven bail-out for non-responsive participants; Gambit
  for extensive-form game generation via frontier exploration;
  SMT-LIB for reachability queries; Scribble for the session-typed
  protocol description; a sketch Bitcoin/Lightning move transcript
  for the two-party fragment (a textual artifact, not deployable
  bytecode); a Coq/Lean ``Diamond DAG'' witness-record encoding
  that captures the per-step constraint bundle as a dependent
  type.

\item \emph{An operational handler layer}
  (\S\ref{sec:liveness}) for liveness and griefing: every action
  can be guarded by a handler that fires if the actor does not
  respond within a deadline.  The handler is a strategic choice
  available to the actor (``quit''), realised by the Solidity
  backend through a timeout-and-bail mechanism, and exposed to the
  game-theoretic backends as an explicit move.

\item \emph{A bridge to a Lean~4 formalization}
  (\S\ref{sec:vegascore}) named VegasCore, which formalizes a core
  calculus over the same EventGraph carrier.  Vegas and VegasCore
  are independent artifacts: Vegas is a compiler tool, VegasCore is
  a self-contained semantic construction.  Where they describe the
  same object they use the same name; where they diverge, the
  divergence is recorded.

\item \emph{An example library} (\S\ref{sec:evaluation}):
  thirty-nine Vegas specifications covering classical mechanisms
  (Vickrey auction, Prisoners' Dilemma, Centipede), Ethereum
  patterns (escrow, micropayment channel, governance voting), and
  adversarial constructions (front-running setups, hidden-reserve
  auctions).  Each example is compiled to every applicable
  backend.
\end{enumerate}

The report follows the order of the compiler.
\S\ref{sec:language} introduces the surface language and walks
through one example end-to-end.  \S\ref{sec:eventgraph} fixes the
pipeline and defines the EventGraph and its invariants.  \S\ref{sec:commit-reveal} describes the
automatic commit-reveal pass.  \S\S\ref{sec:solidity}--\ref{sec:other-backends}
cover the backends.  \S\ref{sec:liveness} describes the
liveness/quit-handler layer.  \S\ref{sec:vegascore} maps the
construction onto the companion Lean formalization.
\S\ref{sec:evaluation} reports on the example suite.
\S\S\ref{sec:related}--\ref{sec:discussion} close with related work
and a frank discussion of current limitations and natural
extensions.
